% Authored LaTeX mathematics notes.
% !TeX program = xelatex
\documentclass{../toolchain/latex/qlnotes-native}
\title{Foundations of Mathematics}
\subtitle{数学逻辑与 ZFC 集合论}
\author{Qiulin Fan}
\date{2026}
\begin{document}
\frontmatter
\maketitle
\tableofcontents
\mainmatter
\chapter{语言、对象与基础}

这门课为之后的数学与计算机课程建立共同的语言：一个表达式是否合式，
它在什么解释下为真，一个结论是否能由指定假设推出，以及所讨论的对象怎样由集合构造。
本轮先学习经典数理逻辑与 ZFC；类型论、范畴论留作之后的独立专题。

\section{两本书的阅读路线}
逻辑部分以 Enderton 的 \emph{A Mathematical Introduction to Logic} 为主：
先读 Chapter Zero 的集合预备知识，再读 Chapter One 的 sentential logic
与 Chapter Two 的 first-order logic。Chapter Three 的 undecidability
与 Chapter Four 的 second-order logic 是后续阅读。

集合论部分以 \emph{Elements of Set Theory} 为主：Chapters 1--2 建立集合观点、
公理和运算；Chapter 3 讨论关系与函数；Chapters 4--5 构造自然数和实数；
Chapter 6 讨论基数与选择公理；Chapters 7--8 讨论序、序数和超限构造。
Chapter 9 的特殊专题留到主线建立后再读。

这些笔记按概念重新组织，并记录初步解释、证明和阅读任务。
书中的完整证明仍需逐节 review；下面标为「证明路线」的内容保留这一区别。

\section{对象语言与元语言}
\begin{definition}{Object language and metalanguage}
\textbf{对象语言}是被研究的形式语言；\textbf{元语言}是我们用来描述该语言、
其解释和推理规则的语言。例如研究集合论时，\(x\in y\) 是对象语言的原子公式，
「这个公式有两个自由变量」是元语言的断言。
\end{definition}

\(\varphi\) 通常是元语言中代表一个公式的字母，不自动成为对象语言里的变量。
同样，\(\models\) 描述满足关系，\(\vdash\) 描述可推导性；它们不是 ZFC
对象语言的关系符号。我们先借助非形式化数学描述语法和语义，随后在 ZFC 中研究这些构造。

\begin{definition}{Working conventions}
本课采用经典一阶逻辑，带等号，解释的论域非空。集合论采用纯集合的 ZFC，
其唯一非逻辑关系符号是 \(\in\)。集合论语境中的所有个体变量都量化于集合。
自然数采用含 \(0\) 的 \(\omega\)；其他课程若把 \(\mathbb N\) 定为正整数，须另行声明。
\end{definition}

\section{每次 review 要回答的四个问题}
\begin{enumerate}
\item \textbf{对象：}变量表示什么，在哪个集合或结构中取值？
\item \textbf{断言：}量词的顺序、作用域和自由变量是什么？
\item \textbf{推导：}用了哪些假设、规则或已经证明的结论？
\item \textbf{构造：}对象的存在性与唯一性由什么保证，是否使用了选择公理？
\end{enumerate}
定义固定术语或引入可消去的记号；公理规定一个理论的假设；定理需要在这些假设下证明。
一个方便的图像、例子或记号本身不保证对象存在。

\part{Mathematical Logic}
\chapter{命题逻辑：语法与语义}
\begin{qlnote}{阅读参照}Enderton, \emph{Logic}, Chapter One.\end{qlnote}

\section{递归定义公式}
\begin{definition}{Sentential formulas}
给定命题字母集合，取 \(\neg,\to\) 为基本联结词。每个命题字母是公式；
若 \(\varphi,\psi\) 是公式，则 \((\neg\varphi)\) 和 \((\varphi\to\psi)\)
也是公式；除此之外没有公式。括号是语法的一部分。
\end{definition}
这一定义同时给出\textbf{结构归纳}：先证明命题字母满足性质，再证明两个构造步骤保留性质，
便可得到所有公式满足该性质。公式是一棵有限语法树；它的语义通过树递归计算。

\section{真值赋值}
\begin{definition}{Truth assignment}
赋值 \(v\) 把每个命题字母映到 \(\{0,1\}\)，并唯一递归扩张到公式：
\[
 v(\neg\varphi)=1-v(\varphi),\qquad
 v(\varphi\to\psi)=0\ \Longleftrightarrow\
 v(\varphi)=1\text{ 且 }v(\psi)=0.
\]
\end{definition}
\(\varphi\land\psi\)、\(\varphi\lor\psi\) 和 \(\varphi\leftrightarrow\psi\)
可作为缩写引入，也可作为带真值规则的基本联结词。前后必须保持同一种约定。

\begin{definition}{Satisfiability and semantic consequence}
若存在赋值使一组公式 \(\Gamma\) 全为真，则 \(\Gamma\) \textbf{可满足}。
\(\Gamma\models\varphi\) 表示每个使 \(\Gamma\) 全为真的赋值也使 \(\varphi\) 为真。
若 \(\varnothing\models\varphi\)，则 \(\varphi\) 是\textbf{重言式}。
若 \(\models\varphi\leftrightarrow\psi\)，则两公式\textbf{逻辑等价}。
\end{definition}

\begin{theorem}{Modus ponens preserves truth}
\(\{\varphi,\varphi\to\psi\}\models\psi\)。
\end{theorem}
\begin{proof}
取使两个前提为真的赋值。若 \(v(\psi)=0\)，由于 \(v(\varphi)=1\)，
条件公式的真值规则给出 \(v(\varphi\to\psi)=0\)，矛盾。
\end{proof}

\section{充分、必要与逆否}
\(P\to Q\) 中，\(P\) 是 \(Q\) 的充分条件，\(Q\) 是 \(P\) 的必要条件。
它与 \(\neg Q\to\neg P\) 逻辑等价。逆命题 \(Q\to P\) 和否命题
\(\neg P\to\neg Q\) 互相等价，却一般不由原命题推出。
赋值 \(v(P)=0,v(Q)=1\) 就使原命题真、逆命题假。

\begin{example}[真值表能检查什么]
一个只含有限多个命题字母的公式，可由有限真值表判定是否为重言式。
例如 \(P\lor\neg P\) 在经典语义下总为真。
这项算法检查的是命题联结结构；它不会检查「连续函数一定可导」这样的数学前提是否成立。
\end{example}

\section{练习}
证明 De Morgan 等价式；为肯定后件 \((P\to Q),Q\therefore P\) 找反例；
用结构归纳证明每个公式只含有限多个命题字母。阅读原书时补足唯一可读性与命题紧致性的证明。

\chapter{一阶逻辑：量词、结构与满足}
\begin{qlnote}{阅读参照}Enderton, \emph{Logic}, Chapter Two.\end{qlnote}

\section{语言、项与公式}
一阶语言的签名规定常元、函数符号和关系符号，并规定后两者的元数。
变量和常元是项；\(f(t_1,\ldots,t_n)\) 是由函数符号构成的项。
原子公式是 \(t=u\) 或 \(R(t_1,\ldots,t_n)\)；再用联结词和量词构造其他公式。
项表示对象，公式表达断言；函数符号和关系符号不能混用。

\begin{definition}{Free variables and sentences}
变量的某次出现若处于同名量词的约束范围内，便是约束出现；否则是自由出现。
没有自由变量的公式称为\textbf{句子}。同一个变量符号可以在不同位置自由或约束地出现。
\end{definition}
例如 \(P(x)\land\forall x\,Q(x)\) 中，左边的 \(x\) 自由，右边的 \(x\) 约束。
\(\forall x(P(x)\to Q(x))\) 则没有自由变量，若 \(P,Q\) 是一元关系符号。

\section{结构与赋值}
\begin{definition}{Structure}
\(\mathcal M\) 给出非空论域 \(M\)，每个常元的元素解释、每个 \(n\) 元函数符号的
全函数 \(M^n\to M\)，以及每个 \(n\) 元关系符号的子集 \(M^n\)。
等号解释为论域上的实际相等。赋值 \(s\) 把变量映到 \(M\)。
\end{definition}
项的值 \(t^{\mathcal M}[s]\) 递归定义，公式的满足关系也递归定义。
量词的两条规则是
\[
\begin{aligned}
\mathcal M,s\models\forall x\,\varphi
 &\Longleftrightarrow \text{每个 }a\in M\text{ 都有 }\mathcal M,s[x\mapsto a]\models\varphi,\\
\mathcal M,s\models\exists x\,\varphi
 &\Longleftrightarrow \text{存在 }a\in M\text{ 使 }\mathcal M,s[x\mapsto a]\models\varphi.
\end{aligned}
\]
句子的真值不依赖赋值。写 \(\Gamma\models\varphi\) 时，要遍历所有结构和赋值；
只在某个固定结构中为真，是更具体的断言。

\section{量词顺序与变量捕获}
\begin{example}[见证是否可以随输入变化]
在实数的通常序结构中，\(\forall x\exists y\,(x<y)\) 为真：可取 \(y=x+1\)。
\(\exists y\forall x\,(x<y)\) 为假：给定 \(y\)，取 \(x=y\)。
前者允许见证依赖于 \(x\)，后者要求同一个见证适用于所有 \(x\)。
\end{example}

\(\varphi[t/x]\) 表示把自由出现的 \(x\) 替换为项 \(t\)，且不得捕获 \(t\) 的自由变量。
例如直接把 \(y\) 代入 \(\exists y(x\ne y)\) 中的 \(x\)，会错误得到
\(\exists y(y\ne y)\)。正确做法是先把约束变量改名为新变量 \(z\)，再得到
\(\exists z(y\ne z)\)。这种改名不改变原公式的含义。

\section{把分析中的条件展开}
一致连续性的量词顺序为
\[
\forall\varepsilon>0\ \exists\delta>0\ \forall x,y\in A\quad
 |x-y|<\delta\ \Longrightarrow\ |f(x)-f(y)|<\varepsilon.
\]
这里 \(\delta\) 可以依赖 \(\varepsilon\)，却必须同时适用于所有 \(x,y\in A\)。
复习普通连续性时，把 \(x\) 放在 \(\exists\delta\) 之前，便能看清两者的区别。

\chapter{推导、可靠性、完备性与紧致性}
\begin{qlnote}{阅读参照}Enderton, \emph{Logic}, Chapters One--Two。
下面用自然演绎展示推理；阅读原书时另行对照其公理式演算。\end{qlnote}

\section{有限推导与假设的解除}
\(\Gamma\vdash\varphi\) 表示存在一个有限的形式推导，从 \(\Gamma\) 中的假设推出
\(\varphi\)。推导的每一步须由规定的规则允许；推导树的叶子记录尚未解除的假设。
以下规则中的 \(\neg\varphi\) 可理解为 \(\varphi\to\bot\)。

\begin{align*}
\frac{\Gamma\vdash\varphi\quad\Gamma\vdash\psi}{\Gamma\vdash\varphi\land\psi}
&\qquad
\frac{\Gamma\vdash\varphi\land\psi}{\Gamma\vdash\varphi}
\qquad
\frac{\Gamma\vdash\varphi\land\psi}{\Gamma\vdash\psi},\\
\frac{\Gamma,\varphi\vdash\psi}{\Gamma\vdash\varphi\to\psi}
&\qquad
\frac{\Gamma\vdash\varphi\to\psi\quad\Gamma\vdash\varphi}{\Gamma\vdash\psi},\\
\frac{\Gamma\vdash\varphi}{\Gamma\vdash\varphi\lor\psi}
&\qquad
\frac{\Gamma\vdash\psi}{\Gamma\vdash\varphi\lor\psi},\\
\frac{\Gamma\vdash\varphi\lor\psi\quad\Gamma,\varphi\vdash\chi\quad\Gamma,\psi\vdash\chi}
 {\Gamma\vdash\chi}
&\qquad
\frac{\Gamma\vdash\bot}{\Gamma\vdash\chi}.
\end{align*}
假设规则允许从 \(\Gamma,\varphi\) 推出 \(\varphi\)；背景假设可加入而不必使用。
条件引入和分情况证明会解除相应的临时假设。
经典反证规则还允许从 \(\Gamma,\neg\varphi\vdash\bot\) 推出 \(\Gamma\vdash\varphi\)。

\section{量词规则的附加条件}
全称消去允许由 \(\forall x\varphi\) 推出 \(\varphi[t/x]\)，前提是代入不捕获变量。
全称引入允许由 \(\Gamma\vdash\varphi\) 推出 \(\Gamma\vdash\forall x\varphi\)，
但 \(x\) 不能在 \(\Gamma\) 的未解除假设中自由出现。
\(P(x)\vdash P(x)\) 因而不允许直接推出 \(P(x)\vdash\forall xP(x)\)。

存在引入由 \(\varphi[t/x]\) 得到 \(\exists x\varphi\)。存在消去先取一个新变量
\(y\) 表示见证，以 \(\varphi[y/x]\) 为临时假设推出 \(\psi\)，再解除该假设。
\(y\) 不得在其他未解除假设、结论 \(\psi\) 或原存在公式中自由出现；代入仍须无捕获。
等号规则包括 \(t=t\)，以及在无捕获代入的公式中以相等项替换。

\section{两个不同的元定理}
\begin{theorem}{Soundness}
对上述经典演算，若 \(\Gamma\vdash\varphi\)，则 \(\Gamma\models\varphi\)。
\end{theorem}
\begin{proof}
对有限推导树归纳。假设在使背景假设全为真的解释中为真；联结词规则由真值规则验证。
全称引入时改变 \(x\) 的赋值不会改变背景假设的真值，这正是其附加条件的作用。
存在消去把新变量赋为一个见证，其附加条件保证背景与结论不依赖这个变量。
等号替换保持对象相同；反证规则使用经典二值语义。故根结论在每个这样的解释中为真。
\end{proof}

\begin{theorem}{Completeness of first-order logic}
对经典一阶逻辑的完备演算，\(\Gamma\models\varphi\) 蕴含 \(\Gamma\vdash\varphi\)。
\end{theorem}
\begin{qlnote}{证明路线}
对句子理论，设 \(\Gamma\nvdash\varphi\)。先证明 \(\Gamma\cup\{\neg\varphi\}\)
一致，再扩充语言加入存在见证，扩展为完备且一致的 Henkin 理论，最后以闭项的等价类
构造模型并证明 truth lemma。这样得到使 \(\Gamma\) 真而 \(\varphi\) 假的模型。
这里记录的是构造路线；见证扩充、一致性保留和 truth lemma 的完整证明须随原书逐步补足。
\end{qlnote}

\begin{theorem}{Compactness}
一阶句子集合 \(\Gamma\) 可满足，当且仅当其每个有限子集可满足。
\end{theorem}
\begin{proof}
一个模型当然满足每个有限子集。反之，若 \(\Gamma\) 不可满足，则
\(\Gamma\models\bot\)。由完备性有 \(\Gamma\vdash\bot\)，而这个有限推导只使用
有限多个来自 \(\Gamma\) 的假设。可靠性说明这些假设组成的有限子集不可满足，矛盾。
\end{proof}

逻辑完备性把语义后承与形式可推导性对应起来；它不会自动证明某个指定理论一致，
也不表示 ZFC 能决定其语言中的每个句子。不可判定性和不完备性留到 \emph{Logic}
Chapter Three 单独学习。

\part{ZFC Set Theory}
\chapter{ZFC：语言、公理与公理模式}
\begin{qlnote}{阅读参照}Enderton, \emph{Set Theory}, Chapters 1--2、6--7 及书末公理表。
这里把不同阶段引入的公理先集中展示，再分章解释其用途。\end{qlnote}

\section{语言与缩写}
纯集合论语言只有非逻辑二元关系符号 \(\in\)，另带逻辑等号。
例如 \(A\subseteq B\) 缩写为 \(\forall x(x\in A\to x\in B)\)。
\(\forall x\in A\,\varphi\) 缩写为 \(\forall x(x\in A\to\varphi)\)，
\(\exists x\in A\,\varphi\) 缩写为 \(\exists x(x\in A\land\varphi)\)。
下文的空集、无序对、并集、幂集、函数等记号，都须由存在唯一性建立后作为定义性缩写使用。

\textbf{公理模式}给每个符合条件的公式一个公理实例。分离与替代是模式，
因而这张有限的名称表不意味着 ZFC 只有有限条一阶公理。
下列各式均作全称闭包；\(\bar p\) 表示公式的有限参数列。
模式中的 \(\varphi\) 是元语言占位符，不能包含正在新建的集合变量 \(B\) 的自由出现。

\section{九组公理}
\begin{axiom}{Extensionality：外延}
\[
 \forall A\forall B\bigl[\forall x(x\in A\leftrightarrow x\in B)\to A=B\bigr].
\]
集合由其成员决定。证明两个集合相等通常从双向包含开始。
\end{axiom}

\begin{axiom}{Pairing：无序对}
\[
 \forall a\forall b\exists C\forall x\bigl(x\in C\leftrightarrow(x=a\lor x=b)\bigr).
\]
由外延性，此集合唯一，记为 \(\{a,b\}\)；取 \(a=b\) 得单点集。
\end{axiom}

\begin{axiom}{Union：并集}
\[
 \forall A\exists U\forall x\bigl(x\in U\leftrightarrow\exists B(B\in A\land x\in B)\bigr).
\]
记为 \(U=\bigcup A\)。二元并集是 \(a\cup b=\bigcup\{a,b\}\)。
\end{axiom}

\begin{axiom}{Power Set：幂集}
\[
 \forall A\exists P\forall x\bigl(x\in P\leftrightarrow x\subseteq A\bigr).
\]
记为 \(P=\mathcal P(A)\)。它保证所有 \(A\) 的子集组成一个集合。
\end{axiom}

\begin{axiom}{Separation schema：分离}
对于集合论公式 \(\varphi(x,\bar p)\)，有公理实例
\[
 \forall\bar p\forall A\exists B\forall x
 \bigl(x\in B\leftrightarrow(x\in A\land\varphi(x,\bar p))\bigr).
\]
它从一个已有集合 \(A\) 筛选成员；\(B=\{x\in A:\varphi(x,\bar p)\}\)。
\end{axiom}

\begin{axiom}{Replacement schema：替代}
对于公式 \(\varphi(x,y,\bar p)\)，有公理实例
\[
\forall\bar p\forall A\left[
 \bigl(\forall x\in A\ \exists!y\,\varphi(x,y,\bar p)\bigr)
 \to\exists B\forall y\bigl(y\in B\leftrightarrow\exists x\in A\,\varphi(x,y,\bar p)\bigr)
\right].
\]
一个在集合域上唯一确定输出的可定义对应，其值域也是集合。
输出没有被预先限制在某个给定集合中。
\end{axiom}

\begin{axiom}{Infinity：无穷}
存在集合 \(I\)，满足
\[
 \varnothing\in I\quad\text{且}\quad\forall x\in I\ (x\cup\{x\}\in I).
\]
这样的 \(I\) 称为归纳集合。空集记号的存在性由下面的推论保证。
\end{axiom}

\begin{axiom}{Foundation：正则／基础}
\[
 \forall A\bigl[A\ne\varnothing\to\exists x\in A\ (x\cap A=\varnothing)\bigr].
\]
每个非空集合具有一个与它无共同成员的成员。交集记号可由分离定义。
\end{axiom}

\begin{axiom}{Choice：选择}
对于每个成员都非空的集合 \(A\)，存在函数 \(f\)，使
\[
 \operatorname{dom}(f)=A\quad\text{且}\quad\forall a\in A\ (f(a)\in a).
\]
函数与定义域的符号将在下一章展开为集合论定义。这是纯一阶语言句子的缩写。
\end{axiom}
本笔记把外延、对、并、幂、分离、替代、无穷、基础合称 ZF，再加入选择得到 ZFC。
有些公理表单列空集公理，或使用更少的独立公理组；计数依所选表达方式而异。

\section{空集、类与 Russell 问题}
\begin{theorem}{空集存在且唯一}
在本课非空论域的一阶逻辑约定下，由分离和外延可以推出空集存在且唯一。
\end{theorem}
\begin{proof}
取一个集合 \(A\)，以条件 \(x\ne x\) 分离，得到没有成员的集合。
任意两个没有成员的集合有相同成员，由外延性相等。
\end{proof}

\begin{theorem}{不存在包含所有集合的集合}
不存在集合 \(U\) 满足 \(\forall x(x\in U)\)。
\end{theorem}
\begin{proof}
若有这样的 \(U\)，以分离得到 \(R=\{x\in U:x\notin x\}\)。
由于 \(R\in U\)，有 \(R\in R\leftrightarrow R\notin R\)，矛盾。
\end{proof}
\(\{x:\varphi(x)\}\) 可在元语言中描述一个\textbf{可定义类}；
只有证明某个集合恰有这些成员后，才能把它当作集合使用。
\(V\)、所有序数组成的类等记号，不是 ZFC 中额外的个体变量或集合常元。

\begin{theorem}{没有集合属于自己}
基础公理蕴含 \(\forall x(x\notin x)\)。
\end{theorem}
\begin{proof}
若 \(x\in x\)，对非空单点集 \(\{x\}\) 应用基础公理。
其唯一成员是 \(x\)，但 \(x\cap\{x\}\) 含有 \(x\)，矛盾。
\end{proof}

\chapter{关系、函数与索引族}
\begin{qlnote}{阅读参照}Enderton, \emph{Set Theory}, Chapters 2--3.\end{qlnote}

\section{有序对与笛卡尔积}
\begin{definition}{Kuratowski ordered pair}
\[
 (a,b):=\{\{a\},\{a,b\}\}.
\]
对公理保证这确实是集合。它须满足 \((a,b)=(c,d)\) 当且仅当 \(a=c,b=d\)。
\end{definition}
\begin{proof}
该对的成员的交集为 \(\{a\}\)，并集为 \(\{a,b\}\)，所以可先恢复 \(a\)。
若并集只有 \(a\)，则 \(b=a\)；否则 \(b\) 是并集中唯一不同于 \(a\) 的元素。
因此可从编码唯一恢复两个分量，反方向由定义立即成立。
\end{proof}

\(A\times B=\{(a,b):a\in A,b\in B\}\) 存在：每个这样的对属于
\(\mathcal P(\mathcal P(A\cup B))\)，再由分离筛出所需对。
这说明「有某种描述」之外，还要给出构造所处的已有集合。

\section{关系与函数}
\begin{definition}{Relations and functions}
从 \(A\) 到 \(B\) 的关系是 \(A\times B\) 的一个子集。
函数图像 \(f\subseteq A\times B\) 满足
\[
 \forall a\in A\ \exists!b\in B\ ((a,b)\in f).
\]
定义 \(\operatorname{dom}(f)=A\)，\(\operatorname{ran}(f)=\{b\in B:\exists a\in A\,(a,b)\in f\}\)，
并写 \(f:A\to B\)。每个函数值由唯一性确定。
\end{definition}
按函数图像的集合定义，陪域 \(B\) 是箭头声明的一部分，却不由图像唯一决定。
例如同一个恒等图像可声明为 \(\{0\}\to\{0\}\) 或 \(\{0\}\to\{0,1\}\)；
只有前一个声明是满射。讨论满射、复合与逆函数时，须写出所声明的陪域。

若 \(f:A\to B,g:B\to C\)，则
\[
 g\circ f=\{(a,c)\in A\times C:\exists b\in B\ ((a,b)\in f\land(b,c)\in g)\}.
\]
分离保证这个图像存在，两个原函数的唯一性保证它是函数。
关系的逆总能定义；逆关系成为反向全函数需要原函数为双射。

\section{等价类与商集合}
等价关系须具有自反、对称、传递性。对 \(E\subseteq A\times A\)，定义
\([a]_E=\{x\in A:(a,x)\in E\}\)。这些类属于 \(\mathcal P(A)\)，故
\[
 A/E=\{C\in\mathcal P(A):\exists a\in A\ (C=[a]_E)\}
\]
由分离存在。在商上定义运算，需要证明改变代表元不会改变所得等价类。

\section{索引族与空索引的交}
集合索引族 \((A_i)_{i\in I}\) 可看成定义域为集合 \(I\) 的函数，其函数值都是集合。
先得到值域集合，再作并集，便得到 \(\bigcup_{i\in I}A_i\)。
当 \(I\ne\varnothing\) 时，任选 \(i_0\in I\)，在 \(A_{i_0}\) 内分离出
满足 \(\forall i\in I\ x\in A_i\) 的元素，就得到交集。

若 \(I=\varnothing\)，不加限制的「所有集合的交」会要求一个全体集合的集合。
在给定环境集合 \(X\) 中定义交，则
\(\{x\in X:\forall i\in\varnothing\ x\in A_i\}=X\)，因为全称条件空真。
这也是分析、拓扑和测度论中须明确环境集合的原因。

\chapter{自然数、归纳与递归}
\begin{qlnote}{阅读参照}Enderton, \emph{Set Theory}, Chapter 4.\end{qlnote}

\section{从无穷公理得到 \texorpdfstring{\(\omega\)}{omega}}
定义后继 \(S(x)=x\cup\{x\}\)。无穷公理给出某个归纳集合 \(I\)。
在 \(\mathcal P(I)\) 内以分离选出所有归纳子集，得到非空集合
\(\mathcal J=\{J\subseteq I:J\text{ 是归纳集合}\}\)，因为 \(I\in\mathcal J\)。
定义 \(\omega=\bigcap\mathcal J\)。

\begin{theorem}{最小归纳集合}
\(\omega\) 是归纳集合，而且包含于每个归纳集合。
\end{theorem}
\begin{proof}
每个 \(J\in\mathcal J\) 都含空集，所以交集含空集；后继闭合性同样保留于交集。
若 \(K\) 是任意归纳集合，则 \(I\cap K\) 是归纳子集，属于 \(\mathcal J\)。
故 \(\omega\subseteq I\cap K\subseteq K\)。
\end{proof}
这样避免把「所有归纳集合」未经证明地当成一个集合。
在这个构造中
\[
 0=\varnothing,\quad 1=\{0\},\quad 2=\{0,1\},\quad 3=\{0,1,2\},\quad n+1=S(n).
\]
每个自然数是其所有严格较小自然数组成的集合；这项性质须由归纳建立。

\section{归纳原理}
\begin{theorem}{Induction on \texorpdfstring{\(\omega\)}{omega}}
若 \(A\subseteq\omega\)，且 \(0\in A\)、\(n\in A\Rightarrow S(n)\in A\)，则 \(A=\omega\)。
\end{theorem}
\begin{proof}
条件说明 \(A\) 是归纳集合，由最小性有 \(\omega\subseteq A\)，结合反向包含得到相等。
\end{proof}
对公式 \(\varphi(n,\bar p)\)，先以分离得到满足该公式的自然数集合，再应用上述定理。
因此「对每个公式成立的归纳」是一族实例；它不是在一阶语言中直接量化于公式。

\section{递归定理与运算}
\begin{theorem}{Recursion on \texorpdfstring{\(\omega\)}{omega}}
给定集合 \(A\)、\(a\in A\) 及函数 \(g:A\to A\)，存在唯一函数 \(h:\omega\to A\)，满足
\(h(0)=a\) 和 \(h(S(n))=g(h(n))\)。
\end{theorem}
\begin{qlnote}{证明路线}
先对每个有限初始段归纳构造递归图像，并证明它们在共同定义域上一致。
这些图像是 \(\omega\times A\) 的子集；可在其幂集中分离出全部有限近似图像，
取并得到 \(h\)。归纳保证全域性，兼容性保证单值性；再用归纳证明两个满足条件的函数相同。
逐项补足有限初始段和近似图像的定义，是本章的原书阅读任务。
\end{qlnote}

由递归可定义
\[
 m+0=m,\quad m+S(n)=S(m+n),\qquad
 m\cdot0=0,\quad m\cdot S(n)=m\cdot n+m.
\]
定义之后还须用归纳证明结合律、交换律及相应消去律。
\textbf{归纳}建立性质；\textbf{递归}构造对象。递推式需要存在唯一性定理才能成为有效定义。

\chapter{整数、有理数与实数的构造}
\begin{qlnote}{阅读参照}Enderton, \emph{Set Theory}, Chapter 5.\end{qlnote}

\section{用等价类补入运算}
自然数的算术及消去律建立后，在 \(\omega\times\omega\) 上定义
\[
 (a,b)\sim(c,d)\quad\Longleftrightarrow\quad a+d=b+c.
\]
整数是这些等价类组成的商集合；\([(a,b)]\) 表示形式差 \(a-b\)。定义
\[
\begin{aligned}
 [(a,b)]+[(c,d)]&=[(a+c,b+d)],\\
 [(a,b)]\,[(c,d)]&=[(ac+bd,ad+bc)].
\end{aligned}
\]
证明 \(\sim\) 为等价关系、运算与代表元选择无关，以及环的各项定律，
是这个构造的必要步骤。\(n\mapsto[(n,0)]\) 给出自然数到整数的嵌入。

有理数同样可由 \(\mathbb Z\times\mathbb Z_{>0}\) 上的等价关系
\((a,b)\sim(c,d)\Longleftrightarrow ad=bc\) 构造。
取正分母便于定义序；加法和乘法仍要先证明良定义。
通常写 \(\omega\subseteq\mathbb Z\subseteq\mathbb Q\subseteq\mathbb R\) 时，
须说明已通过相应嵌入作了识别；这些构造的原始集合不自动按字面互相包含。

\section{Dedekind cuts}
\begin{definition}{Dedekind cut}
有理数的子集 \(D\) 是一个割，若：\(D\ne\varnothing\)、\(D\ne\mathbb Q\)；
若 \(q\in D,r<q\)，则 \(r\in D\)；\(D\) 没有最大元素。
全体割是 \(\mathcal P(\mathbb Q)\) 内以分离得到的一个集合，记为 \(\mathbb R\)。
\end{definition}
定义 \(D\le E\) 当且仅当 \(D\subseteq E\)。对任意两个割，若互不包含，
取 \(d\in D\setminus E,e\in E\setminus D\)，比较 \(d,e\) 后用向下闭合便得矛盾。
所以这个序是线性的。有理数 \(q\) 的嵌入为 \(\{r\in\mathbb Q:r<q\}\)。

\begin{theorem}{Least-upper-bound property of cuts}
若 \(\mathcal A\) 是非空、在包含序下有上界的割集合，则
\(\sup\mathcal A=\bigcup\mathcal A\)。
\end{theorem}
\begin{proof}
令 \(U=\bigcup\mathcal A\)。取一个割上界 \(B\)，有 \(U\subseteq B\ne\mathbb Q\)，
而 \(\mathcal A\) 非空保证 \(U\ne\varnothing\)。向下闭合由各割继承。
每个 \(q\in U\) 属于某个 \(D\in\mathcal A\)，而 \(D\) 无最大元素，故有
\(r\in D\subseteq U\) 满足 \(q<r\)。因此 \(U\) 是割。
并集包含每个 \(D\in\mathcal A\)，且包含于每个上界，故是最小上界。
\end{proof}
实数加法可定义为
\(D+E=\{q\in\mathbb Q:\exists d\in D\exists e\in E\ (q<d+e)\}\)。
乘法、加法逆元和域公理的完整验证随原书补足。本节已经证明的是割序的确界性质，
不能只凭这一条性质就跳过有序域结构的构造。

\chapter{基数、可数性与选择公理}
\begin{qlnote}{阅读参照}Enderton, \emph{Set Theory}, Chapter 6.\end{qlnote}

\section{先比较集合，再引入基数记号}
写 \(A\approx B\) 表示存在双射 \(A\to B\)，写 \(A\preccurlyeq B\) 表示存在单射
\(A\to B\)。前者是等势关系；后者不是两个集合之间的包含关系。
本课称 \(A\) 可数，若 \(A\preccurlyeq\omega\)；可数无穷则要求 \(A\approx\omega\)。
阅读别的课程时，要确认其「可数」是否包含有限集合。

\begin{theorem}{Cantor's theorem}
不存在满射 \(f:A\to\mathcal P(A)\)。
\end{theorem}
\begin{proof}
假设 \(f\) 满射。由分离得到 \(D=\{a\in A:a\notin f(a)\}\)。
\(D\in\mathcal P(A)\)，故存在 \(d\in A\) 使 \(f(d)=D\)。
于是 \(d\in D\Longleftrightarrow d\notin f(d)\Longleftrightarrow d\notin D\)，矛盾。
\end{proof}
单射 \(a\mapsto\{a\}\) 又给出 \(A\preccurlyeq\mathcal P(A)\)。上述对角证明没有使用选择公理。

\begin{theorem}{Cantor--Schröder--Bernstein}
若 \(A\preccurlyeq B\) 且 \(B\preccurlyeq A\)，则 \(A\approx B\)。
\end{theorem}
\begin{proof}
取单射 \(f:A\to B,g:B\to A\)，令 \(A_0=A\setminus g[B]\)，
\(A_{n+1}=g[f[A_n]]\)，\(C=\bigcup_{n\in\omega}A_n\)。定义
\[
 h(a)=\begin{cases}f(a),&a\in C,\\g^{-1}(a),&a\notin C.\end{cases}
\]
若 \(a\notin C\)，则 \(a\notin A_0\)，所以逆值存在，且由单射性唯一。
\(g^{-1}\) 把 \(A\setminus C\) 双射到 \(B\setminus f[C]\)：
若其值落在 \(f[C]\)，则 \(a\in g[f[C]]\subseteq C\)，矛盾；
若 \(b\notin f[C]\)，则 \(g(b)\notin A_0\)，也不能属于任何 \(A_{n+1}\)，否则
由 \(g\) 单射得 \(b\in f[A_n]\)。因此两个分支分别双射到不相交且覆盖 \(B\) 的部分。
\end{proof}

\section{选择的内容与依赖}
选择公理把每个非空集合都有成员的逐项存在断言，提升为整个集合索引族上的一个选择函数。
有限族的选择可由有限归纳得到；任意族的选择是额外公理。
存在明确的统一选取规则时，可直接验证该规则给出函数。

在 ZF 上，选择公理、良序定理和 Zorn 引理等价。
Zorn 引理断言：若一个非空偏序集合的每条链都有该集合内的上界，则存在极大元。
极大元未必是最大元；这两种条件要分开。
这些等价性及基数算术的完整证明，是原书阅读任务。
后续线代用 Zorn 引理证明任意向量空间存在基时，必须记录这一步的选择依赖。

\chapter{序、序数与超限构造}
\begin{qlnote}{阅读参照}Enderton, \emph{Set Theory}, Chapters 7--8。
本章先固定定义与后续证明任务。\end{qlnote}

\section{偏序、线序与良序}
非严格偏序 \(\le\) 满足自反、反对称和传递；线序还要求任意两元素可比较。
严格序 \(<\) 的公理使用反自反性；不能把两套约定混在同一个定义里。
一个严格线序是\textbf{良序}，若每个非空子集都存在最小元素。
\(\omega\) 的通常序是良序；\(\mathbb Z\) 的通常序不是，因为负整数集合没有最小元素。

\section{Von Neumann ordinals}
\begin{definition}{Transitive sets and ordinals}
集合 \(A\) 传递，若 \(x\in y\in A\Rightarrow x\in A\)。
一个传递集合 \(\alpha\)，若其成员被 \(\in\) 严格良序，则称为\textbf{序数}。
对序数写 \(\alpha<\beta\) 当且仅当 \(\alpha\in\beta\)。
\end{definition}
有限自然数和 \(\omega\) 是标准例子。序数后继为 \(\alpha+1=\alpha\cup\{\alpha\}\)；
非零且不是后继的序数称为极限序数。
成员关系与子集关系在序数之间有特殊联系，但不能把这种联系推广到任意集合。

\section{超限归纳与递归的角色}
对一个固定序数 \(\alpha\)，若 \(A\subseteq\alpha\) 满足：对每个 \(\beta<\alpha\)，
只要所有 \(\gamma<\beta\) 属于 \(A\)，就有 \(\beta\in A\)，则 \(A=\alpha\)。
否则 \(\alpha\setminus A\) 的最小元素给出矛盾。

超限递归在阶段 \(\beta\) 根据先前整段函数 \(f\restriction\beta\) 定义新值。
它需要精确说明可定义的规则、值的唯一性及集合长度的定义域。
替代公理用于把集合索引的已定义值收集成集合；极限阶段还需证明近似图像兼容。
序数可比较性、每个良序的序型、超限递归及 rank 的完整构造，随原书逐项展开。

\section{序型与大小}
\(\omega+1\) 的序是在自然数之后添加一个最大元素，所以它不与 \(\omega\) 序同构，
尽管两者的底层集合等势。序数描述良序的型；基数描述集合的大小。
在选择公理下，每个集合可良序，其大小可由某个最小的等势序数，即初始序数表示。
集合的全体与所有序数的全体要区别：后者是一个真类，不能未经论证放进幂集或作为函数的集合定义域。

\part{Review of Mathematics and Computer Science}
\chapter{把基础用于后续课程 review}

\section{统一的检查顺序}
先确定对象和环境，再展开关键量词，检查存在唯一性，然后检查每一步推导及公理依赖。
完成后，把仍未证明的命题、例子提供的证据和已经证明的定理分开记录。

\begin{description}
\item[单变量与多元分析] 区分逐点与一致的量词顺序，写出确界存在的非空、有界条件，
  在构造极限对象时说明它属于哪个集合或空间。
\item[线性代数与代数] 写出底域、向量空间及映射的定义域和陪域；
  商上的运算要证明良定义，无限维存在性证明要标出选择依赖。
\item[测度与概率] 区分点、事件和事件族；\(\sigma\)-代数是 \(\mathcal P(X)\) 的子集，
  其成员才是 \(X\) 的子集。随机变量是指定可测空间之间的函数。
\item[拓扑] 明确环境集合、拓扑及子空间结构；区分「某个指定开覆盖」与「任意开覆盖」。
\item[计算机课程] 区分语法表达式、执行语义和关于程序的断言；
  用循环不变量说明归纳步骤，用终止条件区分部分正确性和全正确性。
\end{description}

\section{第一轮练习清单}
\begin{enumerate}
\item 写出 \(\forall x\exists y\,R(x,y)\) 的否定，并构造它与交换量词后的公式真值不同的结构。
\item 对一个含自由变量的公式做无捕获代入；解释错误代入改变了什么。
\item 从分离与外延证明交集存在且唯一，明确空索引族所用的环境集合。
\item 在 Kuratowski 编码中验证 \((a,a)=(b,c)\) 的情形。
\item 用最小归纳集合证明归纳原理，再给出一个有效递归定义。
\item 检查一个商集合上的运算是否与代表元选择无关。
\item 重写 Cantor 对角证明，标出分离的使用位置。
\item 从已有课程中选一个证明，列出假设、规则、所用定理和选择依赖。
\end{enumerate}

每次实际 review 以一小节或一个命题为单位；这些基础概念在真实证明中的用法，
会反过来决定哪些原书证明需要进一步展开。

\backmatter
\chapter{参考书目}
\begin{enumerate}
\item Herbert B. Enderton, \emph{A Mathematical Introduction to Logic},
  second edition, Academic Press, 2001. ISBN 978-0-12-238452-3.
  \url{https://www.sciencedirect.com/book/9780122384523/a-mathematical-introduction-to-logic}
\item Herbert B. Enderton, \emph{Elements of Set Theory}, Academic Press, 1977.
  ISBN 978-0-12-238440-0.
  \url{https://www.sciencedirect.com/book/9780122384400/elements-of-set-theory}
\end{enumerate}
\end{document}
